\chapter{Sequences and Series}\label{ch:b2:series}

The theory of numerical series (Year 1 volume) matures here in three
directions: series with values in Banach spaces, where absolute
convergence does the work; the finer tests for real series (Abel
summation); and \emph{summable families} --- summation freed from
the order of the terms --- with the Fubini theorem for double sums
and the Cauchy product. These tools carry all the function-series
chapters ahead.

\section{Series in normed spaces}

\begin{definition}\label{def:b2:series:def}
For a sequence $(u_n)$ in a normed space $E$, the series $\sum u_n$
converges when its partial sums do; it converges
\emph{absolutely}\index{absolute convergence!in a Banach space} when
$\sum \norm{u_n} < \infty$. In a \emph{Banach} space, absolute
convergence implies convergence (\cref{thm:b2:nvs:absoluteconvergence});
in a non-complete space this may fail
(\cref{exo:b2:series:9}).
\end{definition}

\begin{example}\label{ex:b2:series:neumann}
In $\mathcal{M}_n(K)$ (or $\mathcal{L}_c(E)$, $E$ Banach): for
$\vertiii A < 1$, the \emph{Neumann series} $\sum A^k$ converges
absolutely to $(I - A)^{-1}$ (proved in \cref{exo:b2:nvs:5});
$\sum \frac{A^k}{k!}$ converges absolutely to $\eu^A$ for every $A$
(\cref{ex:b2:nvs:matrixexp}). Operator-valued geometric and
exponential series behave like their scalar models --- the whole
point of the Banach framework.
\end{example}

\section{Abel summation}

\begin{theorem}[Abel's summation and test]\label{thm:b2:series:abel}
(Summation by parts) For scalars $a_n$ and vectors $b_n$, with $B_n
= \sum_{k=0}^{n} b_k$:\index{Abel summation}
\[
\sum_{n=0}^{N} a_n b_n
= a_N B_N - \sum_{n=0}^{N-1} (a_{n+1} - a_n) B_n .
\]
(Abel's test) If $(a_n)$ is a real sequence, decreasing to $0$, and
the partial sums $B_n$ are \emph{bounded} (in a Banach space), then
$\sum a_n b_n$ converges.
\end{theorem}

\begin{proof}
The identity: write $b_n = B_n - B_{n-1}$ and reindex (a discrete
integration by parts). For the test, with $\norm{B_n} \leq M$: the
boundary term $a_N B_N \to 0$; the series $\sum (a_n -
a_{n+1})B_n$ converges absolutely, since
\[
\sum_n \norm{(a_{n+1} - a_n)B_n} \leq M \sum_n (a_n - a_{n+1})
= M a_0 < \infty
\]
(telescoping, $a_n \downarrow 0$). Both pieces of the identity
converge, hence so does $\sum a_n b_n$.
\end{proof}

\begin{example}\label{ex:b2:series:sinn}
$\sum \frac{\sin n}{n}$ converges: $a_n = \frac1n \downarrow 0$ and
$B_n = \sum_{k=1}^{n} \sin k$ is bounded --- indeed $B_n =
\Im\sum_{k \leq n} \eu^{\iu k} = \Im\,\frac{\eu^{\iu}(\eu^{\iu n} -
1)}{\eu^{\iu} - 1}$, of modulus $\leq \frac{2}{\abs{\eu^{\iu} - 1}}$.
It does \emph{not} converge absolutely ($\abs{\sin n} \geq \sin^2 n
= \frac{1 - \cos 2n}{2}$, and $\sum \frac{1 - \cos 2n}{2n}$ diverges
since $\sum \frac{\cos 2n}{n}$ converges by the same Abel test while
$\sum \frac{1}{2n}$ diverges). The alternating series test is the
special case $b_n = (-1)^n$.
\end{example}

\section{Summable families}

\begin{definition}\label{def:b2:series:summable}
Let $I$ be a countable index set. A family $(u_i)_{i \in I}$ of
\emph{nonnegative reals} is \emph{summable}\index{summable family}
when the finite partial sums are bounded; its sum is
\[
\sum_{i \in I} u_i = \sup_{F \subseteq I \text{ finite}} \sum_{i
\in F} u_i \in \intcc{0}{+\infty} .
\]
A family of reals or complexes (or Banach vectors) is summable when
$(\norm{u_i})$ is; its sum is then defined by splitting into
positive/negative (or real/imaginary) parts --- equivalently, as the
common value of $\sum_{n} u_{\sigma(n)}$ over all enumerations
$\sigma$ of $I$ (see below).
\end{definition}

\begin{theorem}[Summability and order]\label{thm:b2:series:rearrangement}
\begin{enumerate}
  \item For nonnegative families, the sum is invariant under any
        enumeration: $\sum_{i} u_i = \sum_{n=0}^{\infty}
        u_{\sigma(n)}$ for every bijection $\sigma \colon \N \to I$.
  \item A real or complex series $\sum u_n$ is \emph{commutatively
        convergent} (every rearrangement converges, with the same
        sum) if and only if it is absolutely convergent.
\end{enumerate}
\end{theorem}

\begin{proof}
(1) Every partial sum $\sum_{n \leq N} u_{\sigma(n)}$ is a finite
partial sum of the family (so $\leq$ the sup); every finite $F$ is
contained in some $\{\sigma(0), \dots, \sigma(N)\}$ (so the sup
$\leq$ the series' limit). The two bounds match.

(2) If $\sum\abs{u_n} < \infty$: for any rearrangement $\sigma$ and
$\varepsilon > 0$, choose $N$ with $\sum_{n > N}\abs{u_n} \leq
\varepsilon$; beyond the rank where $\sigma$ has exhausted
$\intint{0}{N}$, the rearranged partial sums differ from the
original limit by at most $\varepsilon$: same sum. If $\sum
\abs{u_n} = \infty$ but $\sum u_n$ converges (real case; complex
follows coordinatewise): the positive and negative parts both
diverge, and one can rearrange to reach any prescribed limit
--- Riemann's theorem, carried out in \cref{exo:b2:series:5} ---
so commutative convergence fails.
\end{proof}

\begin{theorem}[Fubini for families; Cauchy products]\label{thm:b2:series:fubini}
Let $(u_{m,n})_{(m,n) \in \N^2}$ be a summable double family (i.e.\
$\sup_F \sum_F \abs{u_{m,n}} < \infty$). Then\index{Fubini theorem
for series}
\[
\sum_{(m,n)} u_{m,n}
= \sum_{m=0}^{\infty}\Bigl(\sum_{n=0}^{\infty} u_{m,n}\Bigr)
= \sum_{n=0}^{\infty}\Bigl(\sum_{m=0}^{\infty} u_{m,n}\Bigr)
= \sum_{k=0}^{\infty}\Bigl(\sum_{m+n=k} u_{m,n}\Bigr),
\]
all inner series converging (absolutely). In particular, if $\sum
a_m$ and $\sum b_n$ converge absolutely, their \emph{Cauchy
product}\index{Cauchy product} converges absolutely with
\[
\Bigl(\sum_m a_m\Bigr)\Bigl(\sum_n b_n\Bigr)
= \sum_{k=0}^{\infty} c_k,
\qquad
c_k = \sum_{m=0}^{k} a_m b_{k-m} .
\]
\end{theorem}

\begin{proof}
\emph{Nonnegative case.} Each grouping (by rows, columns, or
diagonals) computes the same supremum: any finite set of pairs is
contained in a finite block of rows (bounding each grouped sum below
by finite partial sums and above by the total), and monotone
convergence of partial sums does the rest --- concretely, for rows:
$\sum_{m \leq M}\sum_{n \leq N} u_{m,n} \leq S$ gives, letting $N
\to \infty$ then $M \to \infty$, $\sum_m \sum_n u_{m,n} \leq S$;
conversely every finite $F$ sits in such a rectangle, so $S \leq
\sum_m\sum_n u_{m,n}$. Diagonals: same two bounds with triangles
instead of rectangles.

\emph{General case.} Split into positive and negative (real and
imaginary) parts, each a summable nonnegative family; the four
groupings agree on each part, hence on the difference; absolute
convergence of the inner series comes from the nonnegative case
applied to $\abs{u_{m,n}}$.

\emph{Cauchy product.} The family $u_{m,n} = a_m b_n$ is summable:
finite partial sums of $\abs{a_mb_n}$ are bounded by
$\bigl(\sum\abs{a_m}\bigr)\bigl(\sum\abs{b_n}\bigr)$. Rows give
$\bigl(\sum a_m\bigr)\bigl(\sum b_n\bigr)$; diagonals give $\sum_k
c_k$.
\end{proof}

\begin{example}[The exponential identity, honestly]\label{ex:b2:series:exp}
For $a, b \in \C$ (or commuting matrices):
\[
\Bigl(\sum_m \frac{a^m}{m!}\Bigr)\Bigl(\sum_n
\frac{b^n}{n!}\Bigr)
= \sum_k \sum_{m+n=k} \frac{a^m b^n}{m!\,n!}
= \sum_k \frac{(a + b)^k}{k!},
\]
by the binomial theorem on each diagonal: $\eu^a \eu^b = \eu^{a+b}$
--- the functional equation of $\exp$ derived from the series alone.
(Commutation is used in the binomial step; for non-commuting
matrices the identity genuinely fails,
\cref{ch:b2:diffeq}.)
\end{example}

\begin{example}[A double-sum evaluation]\label{ex:b2:series:doublesum}
For $s > 1$ real, let $\zeta(s) = \sum_{n\geq1} n^{-s}$. Counting
divisors by double summation --- the family $(m^{-s}n^{-s})$ over
$(m,n) \in (\N^*)^2$ is summable (product of convergent positive
series) --- and grouping by the product $q = mn$:
\[
\zeta(s)^2 = \sum_{m,n} \frac{1}{(mn)^s}
= \sum_{q=1}^{\infty} \frac{d(q)}{q^s},
\]
where $d(q)$ is the number of divisors of $q$. Summable families
turn combinatorics into analysis.
\end{example}

\section{Exercises}

\begin{exercise}[$\star$]\label{exo:b2:series:1}
Nature of: $\sum \dfrac{\cos n}{n}$; $\;\sum \dfrac{(-1)^n}{\ln
n}$; $\;\sum \dfrac{(-1)^n}{n^{3/4} + \cos n}$ \emph{(expand as in
the Year 1 trap: the alternating test needs monotonicity)}.
\end{exercise}

\begin{exercise}[$\star$]\label{exo:b2:series:2}
Prove that for $\abs z < 1$: $\sum_{n\geq1} n z^{n} =
\dfrac{z}{(1-z)^2}$, via the Cauchy product of $\sum z^n$ with
itself.
\end{exercise}

\begin{exercise}[$\star\star$]\label{exo:b2:series:3}
(Kronecker-type lemma) Let $\sum b_n$ be a convergent real series.
Prove, by Abel summation, that $\dfrac{1}{n}\sum_{k=1}^{n} k\,b_k
\to 0$.
\end{exercise}

\begin{exercise}[$\star\star$]\label{exo:b2:series:4}
Study the convergence of $\sum \dfrac{\sin(n\theta)}{n^\alpha}$
($\theta \in \R$, $\alpha > 0$) --- for which $(\theta, \alpha)$ is
it absolutely convergent, semi-convergent, divergent?
\end{exercise}

\begin{exercise}[$\star\star\star$]\label{exo:b2:series:5}
(Riemann rearrangement) Let $\sum u_n$ be a convergent but not
absolutely convergent real series, and $\ell \in \R$. Prove that
some rearrangement of $\sum u_n$ converges to $\ell$.
\emph{(Show both subseries of positive and negative terms diverge;
then greedily alternate: take positive terms until exceeding
$\ell$, then negative until dropping below, and so on; the terms
tend to $0$, forcing convergence to $\ell$.)}
\end{exercise}

\begin{exercise}[$\star\star$]\label{exo:b2:series:6}
Prove that the family
$\Bigl(\dfrac{x^{m+n}}{m!\,n!}\Bigr)_{(m,n)\in\N^2}$ is summable
for every $x \in \R$, and re-derive the identity $(\eu^x)^2 =
\eu^{2x}$ by grouping the double sum along the diagonals $m + n =
k$.
\end{exercise}

\begin{exercise}[$\star\star$]\label{exo:b2:series:7}
Prove that the family $\bigl(\frac{1}{m^2 n^2}\bigr)_{m,n \geq 1}$
is summable, and that grouping by $\gcd$: with $q = \gcd(m,n)$,
\[
\zeta(2)^2 = \sum_{q\geq1} \frac{1}{q^4}
\sum_{\substack{a,b \geq 1\\ \gcd(a,b)=1}} \frac{1}{a^2b^2}
= \zeta(4) \cdot S,
\]
where $S = \sum_{\gcd(a,b)=1} \frac{1}{a^2b^2}$: deduce $S =
\zeta(2)^2/\zeta(4)$. \emph{(Every pair $(m,n)$ writes uniquely
$(qa, qb)$ with $\gcd(a,b) = 1$.)}
\end{exercise}

\begin{exercise}[$\star\star\star$]\label{exo:b2:series:8}
(Abel's theorem on products, light version) Suppose $\sum a_n$
converges absolutely and $\sum b_n$ converges. Prove that their
Cauchy product $\sum c_n$ converges, with $\sum c_n = (\sum
a_n)(\sum b_n)$. \emph{(Write $C_N = \sum_{k\leq N} c_k = \sum_n
a_n B_{N-n}$ with $B$ the partial sums of $b$; split according to
$n \leq N/2$ or not, using boundedness of $(B_m)$ and the absolute
tail of $(a_n)$.)}
\end{exercise}

\begin{exercise}[$\star\star\star$]\label{exo:b2:series:9}
In the (non-complete) space $E$ of eventually-zero real sequences
with the sup norm, exhibit an absolutely convergent series that does
not converge in $E$. \emph{(Try $u_n = 2^{-n} e_n$ with $(e_n)$ the
canonical sequences.)}
\end{exercise}
